\documentclass[10.5pt,a4paper]{article}
% shared preamble for the English and Bulgarian papers (XeTeX; built with tectonic)
\usepackage[a4paper,left=26mm,right=26mm,top=26mm,bottom=28mm]{geometry}
\usepackage{amsmath}
\usepackage{fontspec}
\setmainfont{Palatino}
\setmonofont{Menlo}[Scale=MatchLowercase]
\usepackage{polyglossia}
\newfontfamily\cyrillicfont{Palatino}
\newfontfamily\cyrillicfonttt{Menlo}[Scale=MatchLowercase]
\usepackage{microtype}
\usepackage{booktabs,array,tabularx}
\usepackage{enumitem}
\usepackage{titlesec}
\usepackage{fancyhdr}
\usepackage{natbib}
\usepackage[hidelinks]{hyperref}
\usepackage{xurl}
\usepackage[font=small,labelfont=bf,labelsep=period,justification=raggedright,singlelinecheck=false,skip=4pt]{caption}

\usepackage{pgfplots}
\usepgfplotslibrary{groupplots}
\pgfplotsset{compat=1.17}
\definecolor{refcolor}{HTML}{0072B2}
\definecolor{varcolor}{HTML}{D55E00}
\definecolor{repcolor}{HTML}{009E73}

\setlength{\parindent}{0pt}
\setlength{\emergencystretch}{2em}
\setlength{\parskip}{5pt plus 1pt}
\linespread{1.07}
\bibpunct{(}{)}{;}{a}{}{,}
\setlength{\bibsep}{1.5pt}
\renewcommand{\bibfont}{\small}
\setlist{nosep,leftmargin=1.4em,itemsep=2pt}
\titleformat{\section}{\large\bfseries}{\thesection}{0.6em}{}
\titleformat{\subsection}{\normalsize\bfseries}{\thesubsection}{0.6em}{}
\titlespacing*{\section}{0pt}{14pt}{5pt}
\titlespacing*{\subsection}{0pt}{9pt}{3pt}
\pagestyle{plain}
\renewcommand{\arraystretch}{1.15}
\newcolumntype{L}[1]{>{\raggedright\arraybackslash}p{#1}}
\newcommand{\code}[1]{\texttt{\small #1}}

\newcommand{\titleblock}[3]{%
  \begin{center}
    {\Large\bfseries #1\par}\vspace{8pt}
    {\large #2\par}\vspace{2pt}
    {\small #3\par}
  \end{center}\vspace{4pt}}
\newenvironment{abstractbox}[1]{\begin{center}\bfseries\small #1\end{center}\vspace{-6pt}\begin{quote}\small}{\end{quote}}

% worked-example figure. arguments: time label, reference, edit inside the network, passive reporter, occupancy title.
% data: figure-patterning.csv, written from the Studio simulation of GAACACCCA (GLI2), signal 0.2 for 20 model units.
\newcommand{\patterningplot}[5]{%
\begin{tikzpicture}
\begin{groupplot}[group style={group size=3 by 2,horizontal sep=15mm,vertical sep=17mm},
  width=44mm,height=30mm,xmin=0,xmax=20,xtick={0,10,20},ymin=0,scaled y ticks=false,
  yticklabel style={/pgf/number format/fixed,/pgf/number format/precision=3},
  tick label style={font=\scriptsize},label style={font=\scriptsize},title style={font=\footnotesize\bfseries},
  every axis plot/.append style={thick}]
\nextgroupplot[title={Pax6}]
\addplot[refcolor] table[x=t,y=refPax6,col sep=comma]{figure-patterning.csv};
\addplot[varcolor,dashed] table[x=t,y=varPax6,col sep=comma]{figure-patterning.csv};
\nextgroupplot[title={Olig2}]
\addplot[refcolor] table[x=t,y=refOlig2,col sep=comma]{figure-patterning.csv};
\addplot[varcolor,dashed] table[x=t,y=varOlig2,col sep=comma]{figure-patterning.csv};
\nextgroupplot[title={Nkx2.2}]
\addplot[refcolor] table[x=t,y=refNkxTwoTwo,col sep=comma]{figure-patterning.csv};
\addplot[varcolor,dashed] table[x=t,y=varNkxTwoTwo,col sep=comma]{figure-patterning.csv};
\nextgroupplot[title={Irx3},xlabel={#1}]
\addplot[refcolor] table[x=t,y=refIrx3,col sep=comma]{figure-patterning.csv};
\addplot[varcolor,dashed] table[x=t,y=varIrx3,col sep=comma]{figure-patterning.csv};
\nextgroupplot[title={#5},xlabel={#1},ymax=1]
\addplot[refcolor] table[x=t,y=occRef,col sep=comma]{figure-patterning.csv};
\addplot[repcolor,densely dotted] table[x=t,y=occRep,col sep=comma]{figure-patterning.csv};
\addplot[varcolor,dashed] table[x=t,y=occVar,col sep=comma]{figure-patterning.csv};
\end{groupplot}
\node[anchor=north west,align=left,font=\scriptsize,inner sep=0pt] at ([xshift=12mm,yshift=-1mm]group c2r2.north east) {%
\tikz[baseline=-0.5ex]\draw[thick,refcolor](0,0)--(7mm,0);~#2\\[4pt]
\tikz[baseline=-0.5ex]\draw[thick,varcolor,dashed](0,0)--(7mm,0);~#3\\[4pt]
\tikz[baseline=-0.5ex]\draw[thick,repcolor,densely dotted](0,0)--(7mm,0);~#4};
\end{tikzpicture}}

\usetikzlibrary{arrows.meta}

% predictor pipeline. arguments: edit pair, window, 5-mer block, motif block, linear score, panel.
\newcommand{\pipelinefig}[6]{%
\begin{tikzpicture}[box/.style={draw,rounded corners=1.5pt,align=center,font=\scriptsize,inner sep=3pt,text width=24mm,minimum height=13mm},
  arr/.style={-{Stealth[length=1.8mm]},thick}]
\node[box] (a) at (0,0) {#1};
\node[box] (b) at (30mm,0) {#2};
\node[box] (c) at (60mm,10mm) {#3};
\node[box] (d) at (60mm,-10mm) {#4};
\node[box] (e) at (90mm,0) {#5};
\node[box] (f) at (120mm,0) {#6};
\draw[arr] (a)--(b);
\draw[arr] (b.east)--(c.west);
\draw[arr] (b.east)--(d.west);
\draw[arr] (c.east)--(e.west);
\draw[arr] (d.east)--(e.west);
\draw[arr] (e)--(f);
\end{tikzpicture}}

% internal signed utility by panel size (build/neural-predictor-reviewed-20261005/siraj-huber-joint-20261019/RESULTS.md).
% arguments: k5, joint-linear, Huber-linear, x label, y label.
\newcommand{\utilityplot}[5]{%
\begin{tikzpicture}
\begin{axis}[ybar,bar width=7pt,width=88mm,height=42mm,symbolic x coords={10,50,200},xtick=data,enlarge x limits=0.25,
  ymin=0,ymax=2.5,ylabel={#5},xlabel={#4},tick label style={font=\scriptsize},label style={font=\scriptsize},
  legend style={font=\scriptsize,at={(0.97,0.97)},anchor=north east,draw=none,fill=none},legend cell align=left,area legend]
\addplot[fill=black!35,draw=none] coordinates {(10,0.857) (50,0.607) (200,0.509)};
\addplot[fill=refcolor,draw=none] coordinates {(10,1.601) (50,0.954) (200,0.626)};
\addplot[fill=varcolor,draw=none] coordinates {(10,1.483) (50,0.986) (200,0.634)};
\legend{#1,#2,#3}
\end{axis}
\end{tikzpicture}}

% Scan of the 200-base example, from figure-scan.csv (docs/paper/figure-data.py). arguments: x label, colour-bar label, example-edit label.
\newcommand{\scanplot}[3]{%
\begin{tikzpicture}
\begin{axis}[width=124mm,height=24mm,scale only axis,axis on top,clip=false,y dir=reverse,xmin=0.5,xmax=200.5,ymin=-0.5,ymax=3.5,
  xtick={1,50,100,150,200},ytick={0,1,2,3},yticklabels={A,C,G,T},xlabel={#1},ylabel={},
  tick label style={font=\scriptsize},label style={font=\scriptsize},
  colormap={eff}{color(0cm)=(refcolor); color(1cm)=(white); color(2cm)=(varcolor)},
  point meta min=-0.6,point meta max=0.6,unbounded coords=jump,
  colorbar,colorbar style={width=2.5mm,ytick={-0.6,0,0.6},yticklabel style={font=\scriptsize,/pgf/number format/fixed,/pgf/number format/precision=1},
    ylabel={#2},ylabel style={font=\scriptsize}}]
\addplot[matrix plot*,mesh/cols=200,point meta=explicit] table[x=pos,y=altidx,meta=effect,col sep=comma]{figure-scan.csv};
\draw[-{Stealth[length=1.6mm]},thick] (axis cs:100,-1.9) -- (axis cs:100,-0.65);
\node[font=\scriptsize,anchor=south] at (axis cs:100,-1.9) {#3};
\end{axis}
\end{tikzpicture}}

% per-study correlation of the primary five studies, from figure-studies.csv. arguments: x label, Huber-linear, S2, k5.
\newcommand{\correlationplot}[4]{%
\begin{tikzpicture}
\begin{axis}[width=58mm,height=44mm,scale only axis,y dir=reverse,ymin=-0.6,ymax=4.6,ytick={0,1,2,3,4},yticklabels={Lee,McAfee,Matoba,Guo,Weiss},
  xmin=-0.05,xmax=0.15,xtick={-0.05,0,0.05,0.10,0.15},scaled x ticks=false,xticklabel style={/pgf/number format/fixed,/pgf/number format/precision=2},xlabel={#1},tick label style={font=\scriptsize},label style={font=\scriptsize},
  xmajorgrids,major grid style={black!12},extra x ticks={0},extra x tick style={grid=major,major grid style={black!55}},extra x tick labels={},
  legend style={font=\scriptsize,at={(0.5,1.04)},anchor=south,draw=none,legend columns=-1,/tikz/every even column/.append style={column sep=3mm}},legend cell align=left]
\addplot[only marks,mark=*,mark size=2.2pt,varcolor,restrict expr to domain={\thisrow{primary}}{1:1}] table[x=huber,y expr=\coordindex,col sep=comma]{figure-studies.csv};
\addplot[only marks,mark=square*,mark size=2pt,refcolor,restrict expr to domain={\thisrow{primary}}{1:1}] table[x=s2,y expr=\coordindex,col sep=comma]{figure-studies.csv};
\addplot[only marks,mark=triangle*,mark size=2.4pt,black!45,restrict expr to domain={\thisrow{primary}}{1:1}] table[x=k5,y expr=\coordindex,col sep=comma]{figure-studies.csv};
\legend{#2,#3,#4}
\end{axis}
\end{tikzpicture}}

% Huber-linear error relative to predicting no effect: internal tuning partition (0.28766713 against 0.30407514) and the five studies
% (figure-studies.csv; the mean bar is 0.10437080 against 0.09529676). arguments: x label, internal, five-study mean.
\newcommand{\excessplot}[3]{%
\begin{tikzpicture}
\begin{axis}[xbar,bar width=6pt,width=52mm,height=44mm,scale only axis,y dir=reverse,ymin=-0.6,ymax=6.6,ytick={0,1,2,3,4,5,6},
  yticklabels={#2,Lee,McAfee,Matoba,Guo,Weiss,#3},xmin=-10,xmax=60,xtick={0,20,40,60},xlabel={#1},
  tick label style={font=\scriptsize},label style={font=\scriptsize},xmajorgrids,major grid style={black!12},
  extra x ticks={0},extra x tick style={grid=major,major grid style={black!55}},extra x tick labels={}]
\addplot[fill=refcolor,draw=none] coordinates {(-5.40,0)};
\addplot[fill=varcolor,draw=none,restrict expr to domain={\thisrow{primary}}{1:1}] table[x=excess,y expr=\coordindex+1,col sep=comma]{figure-studies.csv};
\addplot[fill=black!55,draw=none] coordinates {(9.52,6)};
\end{axis}
\end{tikzpicture}}

\setmainlanguage{bulgarian}
\setotherlanguage{english}
\hypersetup{pdftitle={Към пост-AlphaFold голямото предизвикателство: криптографски проверимо декодиране от некодиращ вариант към функция (V2F) и двукритериално обратно спасяване на виртуална клетка с Epigenesis},pdfauthor={Деян Тодоров}}

\begin{document}

\titleblock{Към пост-AlphaFold голямото предизвикателство: криптографски\\проверимо декодиране от некодиращ вариант към функция (V2F)\\и двукритериално обратно спасяване на виртуална клетка с Epigenesis}{Деян Тодоров}{Научен доклад · 21 ноември 2026 г. (Epigenesis 1.7.0) · English version: \href{epigenesis-v2f-grand-challenge.pdf}{epigenesis-v2f-grand-challenge.pdf}\\Възпроизводим скрипт: \code{examples/developmental-pipeline/grand\_challenge.py} · SHA-256 на доказателството: \code{fa135804ba0b\ldots}}

\begin{abstractbox}{Резюме}
След като предсказването на триизмерната белтъчна структура до голяма степен разреши задачата за 2\% белтъчно-кодиращ геном, централното нерешено голямо предизвикателство в изчислителната биология е \textbf{проблемът „некодиращ вариант към функция“ (Variant-to-Function, V2F) и изграждането на виртуална развиваща се клетка}: предсказването — и обратното инженерно проектиране — на начина, по който промени в 98\% некодираща ДНК се разпространяват през биофизиката на транскрипционните фактори (ТФ), кинетиката на енхансерите в развитието, нелинейната обратна връзка в генните регулаторни мрежи (ГРМ) и триизмерната тъканна морфогенеза. В настоящата работа представяме частично, цялостно и машинно верифицирано решение чрез \textbf{Epigenesis}, което свързва четирите мащаба за 22,14\,с върху централен процесор с междуезиково криптографско възпроизвеждане на Python и на самостоятелен Rust валидатор (\code{epigenesis-validate-rs}). Първо, върху 1895 строго отделени човешки енхансерни редакции в 118 невиждани геномни локуса при неврална индукция (\code{GSE188264}, 0--72\,ч) нашият асиметричен биофизичен модел \textbf{Biophysical Asymmetric 532-TF+ v2} — разделящ разрушаването от създаването на мотиви чрез двумащабна симетрична Болцманова термодинамика, пиков афинитет и контекстно модулиране — достига извънизвадкова корелация на Пиърсън \textbf{$r=0{,}5060$} (95\% бутстрап доверителен интервал по 1000 локуса $[0{,}3938;\, 0{,}6128]$), рангова корелация на Спирман $\rho=0{,}4351$ ($[0{,}3197;\, 0{,}5325]$), устойчива точност във всички седем времеви точки ($r=0{,}4666$--$0{,}4978$), \textbf{30,50\%} намаление на семейно претеглената MSE спрямо нулев ефект ($0{,}250815$ срещу $0{,}360870$), \textbf{96,3\%} точност на посоката в горния децил и \textbf{100,0\%} (Top-25 и Top-50) и \textbf{99,0\%} (Top-100) емпирична лабораторна репресия (точен пермутационен тест $p=0{,}0005$). Второ, количествено доказваме \emph{парадокса на пасивния репортер} в реални човешки енхансери: единични замени, изместващи пасивния репортер с едва $-0{,}012$, претърпяват \textbf{77,8$\times$--78,2$\times$} ендогенно усилване чрез обратна връзка, преобръщат клетъчната съдба и изкривяват триизмерната невритна архитектура в CX3D. Трето, решаваме \textbf{двукритериалната обратна задача за виртуална клетка}: при заключена и непокътната първична мутация Epigenesis синтезира Парето-оптимални дистални еднобазови спасявания, които едновременно възстановяват \textbf{100,00\%} от триизмерната невритна морфология в \code{chr18:56475717-56475888} (\code{165:G>T} + \code{163:T>C}, максимална генна грешка $9{,}89\times10^{-8}$, 100,0\% ансамблово възстановяване) и \textbf{91,13\%} в \code{chr1:175099416-175099587} (\code{149:G>T} + \code{147:A>T}, 95,8\% ансамблово възстановяване), като същевременно намаляват глобалното нецелево 532-ТФ отклонение за 0--72\,ч с \textbf{23,8\%--29,0\%}. И двете решения преминават независима верификация на Python и Rust, експортират RO-Crate~1.1 / COMBINE OMEX архиви и отхвърлят еднобайтова подмяна на входния FASTA. Ясно разграничаваме решеното днес от отворените научни граници за междуизследователски пренос и \emph{in vivo} човешко развитие.

\smallskip\noindent\textbf{Ключови думи:} вариант към функция (V2F), виртуална клетка, некодираща цис-регулаторна граматика, Болцманова термодинамика, Парето-оптимален обратен дизайн на ДНК, триизмерна невронална морфогенеза, GA4GH VRS~2.x, верификация с Rust.
\end{abstractbox}

\section{Пост-AlphaFold голямото предизвикателство}

Нагъването на белтъците определяше изчислителната биология в продължение на половин век; неговото решаване освети около 2\% от човешкия геном, които кодират аминокиселинни вериги. Над 90\% от генетичните варианти, свързани с човешки заболявания — особено при нарушения на невронното развитие — обаче лежат в останалите 98\% некодираща ДНК. За разлика от белтъчно-кодиращите екзони, цис-регулаторната ДНК няма универсална таблица на генетичния код. Нейният функционален смисъл възниква едва през четири свързани физични мащаба:
\begin{enumerate}
\item \textbf{Молекулярна биофизика:} Кооперативно, симетрично по двете вериги свързване на стотици транскрипционни фактори върху припокриващи се ДНК мотиви, при което разрушаването на функционален сайт в активен енхансер е силно асиметрично спрямо създаването на слаб мотив.
\item \textbf{Кинетика на енхансерите в развитието:} Нестационарна регулаторна активност, разгръщаща се във времето на диференциация (напр. 0--72\,ч индукция на човешки плурипотентни стволови клетки към неврални прогенитори).
\item \textbf{Нелинейна обратна връзка в ГРМ (парадоксът на пасивния репортер):} В спрегнати многогенни мрежи слабо смущение на афинитета, което едва измества изолиран репортер, може да прехвърли ендогенния превключвател през седловидно-възлова бифуркация и да обърне клетъчната идентичност.
\item \textbf{Триизмерна тъканна морфогенеза и обратно терапевтично проектиране:} Промените в клетъчната съдба изменят аксонално-дендритното разклоняване и синаптичните контакти; обратно, когато патогенната мутация не може да бъде директно възстановена (поради липса на PAM сайт или структурно сцепление), цис-регулаторното инженерство трябва да синтезира дистална компенсаторна замяна, която възстановява триизмерния фенотип, \emph{без} да поражда нецелева ТФ активност в останалата част на човешкия регулаторен протеом.
\end{enumerate}

Напредъкът по това предизвикателство \textbf{Variant-to-Function (V2F) и виртуална развиваща се клетка} е възпрепятстван не само от точността на моделите, но и от \emph{верифицируемостта}: многомащабните изчисления обикновено живеят в незапечатани скриптове, където незабелязани редакции на входа или смесване на репортерни с ендогенни ефекти компрометират изводите. Тук показваме колко далеч стига \textbf{Epigenesis} в едно детерминистично, напълно възпроизводимо изпълнение (\code{examples/developmental-pipeline/grand\_challenge.py}, 22,14\,с на процесор, канонична SHA-256 сума \code{fa135804ba0b3524\ldots}).

\section{Етап 1: Извънизвадково декодиране на човешка некодираща ДНК (0--72\,ч)}

\subsection{Асиметрична биофизична формулировка (532-TF+ v2)}

Стандартните регресии от последователност към активност третират загубата и придобиването на мотив симетрично ($\Delta\text{активност} \propto w \cdot \Delta\text{оценка}$). Физически обаче човешките енхансери са силно асиметрични: в масово паралелния репортерен експеримент на Kreimer и сътр.\ (2022) за неврална индукция на H1 (\code{GSE188264}, 9241 редакции в 0, 3, 6, 12, 24, 48 и 72\,ч), 69,1\% от пертурбациите потискат активността, а едва 30,9\% я повишават.

Моделът \code{epigenesis-biophysical-temporal-v2} съчетава пет физично обосновани блока признаци върху всички 532 човешки транскрипционни фактора от JASPAR~2024:
\begin{enumerate}
\item \textbf{Двумащабни симетрични по вериги Болцманови лог-суми на състоянията} $\Delta\ln(Z_m + Z_{\mathrm{ns}})$ при $Z_{\mathrm{ns}}\in\{2\times10^{-4},\, 2\times10^{-2}\}$, разделени изрично на асиметрична загуба $\min(0,\Delta\ln Z_m)$ и придобиване $\max(0,\Delta\ln Z_m)$.
\item \textbf{Пиково разрушаване и създаване на сайт} $\min(0,\Delta\max_i s_{m,i})$ и $\max(0,\Delta\max_i s_{m,i})$, отделящо фокалното унищожаване на водещия ТФ сайт от фоновия афинитет.
\item \textbf{Контекстно модулирано разрушаване} $\tanh(\text{ref\_load}_m)\cdot \min(0,\Delta\ln Z_m)$, претеглящо загубата на мотив според реалното му присъствие в референтния енхансер.
\item \textbf{Асиметрични компресирани мотивни потенциали} и \textbf{многомащабни канонични 3-мери и 5-мери}.
\item \textbf{Ортогонално QR времево разлагане}, отделящо средния 0--72\,ч ефект от шест времеви контраста със семейно претеглена гребенова регуляризация, настроена строго върху \code{train}$\to$\code{validation} преди еднократната оценка върху 118-те строго отделени тестови локуса.
\end{enumerate}

\subsection{Строго извънизвадково емпирично доказателство ($n=1895$, 118 локуса)}

Таблица~\ref{tab:v2f} и фигура~\ref{fig:v2f_hours} представят доказателството върху \code{GSE188264}. Върху всичките 1895 невиждани човешки редакции корелацията на Пиърсън достига \textbf{$r=0{,}5060$} (95\% бутстрап ДИ по 1000 семейства $[0{,}3938;\, 0{,}6128]$), а ранговата корелация на Спирман е \textbf{$\rho=0{,}4351$} ($[0{,}3197;\, 0{,}5325]$), с устойчиви стойности във всеки час на диференциация ($r=0{,}4666$--$0{,}4978$, с пикове в 12-ия и 48-ия час). Семейно претеглената целева MSE спада с \textbf{30,50\%} под нулев ефект ($0{,}250815$ срещу $0{,}360870$), надминавайки динуклеотидния ($0{,}319358$), 5-мерния ($0{,}285185$) и симетричния модел на Kreimer v1 ($0{,}252079$).

За реалния лабораторен подбор точността на посоката в горния децил достига \textbf{96,3\%}, а обратният подбор постига \textbf{100,0\%} реална лабораторна репресия при Top-25 (25/25, средно $\Delta=-1{,}3125$, точен 2000-пермутационен тест $p=0{,}0005$), \textbf{100,0\%} при Top-50 (50/50, средно $\Delta=-0{,}9497$, $p=0{,}0005$) и \textbf{99,0\%} при Top-100 (99/100, средно $\Delta=-0{,}8029$, разлика нагоре$-$надолу $+0{,}6530$, $p=0{,}0005$). Сред водещите невиждани попадения моделът автоматично откроява ключови човешки неврорегулаторни ТФ (\code{MA1962.1}: 157 попадения, \code{MA0786.2}: 96, \code{MA0792.1}: 59, \code{MA0841.2}: 11, \code{MA0867.3}: 10).

\begin{table}[tb]
\scriptsize\setlength{\tabcolsep}{4pt}
\caption{Доказателство за Етап 1 върху 1895 строго невиждани човешки редакции при неврална индукция в 118 отделени локуса (\code{GSE188264}, 0--72\,ч).}\label{tab:v2f}
\begin{tabular}{@{}L{41mm}rrrrrr@{}}
\toprule
Модел / лабораторен подбор & Пиърсън $r$ & Спирман $\rho$ & Цел.\ MSE & Горен децил & Репресия \% & Перм.\ $p$ \\
\midrule
Нулев базов модел & 0,0000 & 0,0000 & 0,360870 & 50,0\% & 69,1\% & — \\
Динуклеотиди / k5 гребенова & — & — & 0,285185 & — & — & — \\
Kreimer 0--72\,ч симетричен v1 & — & — & 0,252079 & — & — & — \\
Чист 532-ТФ Болцман (0 $k$-мери) & 0,3439 & — & — & 93,2\% & 95,0\% (Top-100) & 0,000999 \\
\textbf{Biophysical Asymmetric v2} & \textbf{0,5060} & \textbf{0,4351} & \textbf{0,250815} & \textbf{96,3\%} & \textbf{99,0\%} (Top-100) & \textbf{0,0005} \\
\midrule
v2 Top-25 / Top-50 подбор & \multicolumn{6}{@{}L{106mm}@{}}{\raggedright Top-25: \textbf{100,0\%} (25/25, $\bar{\Delta}=-1{,}3125$, $p=0{,}0005$) · Top-50: \textbf{100,0\%} (50/50, $\bar{\Delta}=-0{,}9497$, $p=0{,}0005$)} \\
v2 95\% бутстрап ДИ & \multicolumn{6}{@{}L{106mm}@{}}{\raggedright Пиърсън $r\in [0{,}3938;\, 0{,}6128]$ · Спирман $\rho\in [0{,}3197;\, 0{,}5325]$ · Целеви $r=0{,}4987$, $\rho=0{,}4531$} \\
\bottomrule
\end{tabular}
\end{table}

\begin{figure}[tb]
\centering
\begin{tikzpicture}
\begin{axis}[width=72mm,height=38mm,xlabel={Време на неврална индукция (часове)},ylabel={Извънизвадкова Пиърсън $r$},
  xmin=0,xmax=72,xtick={0,12,24,48,72},ymin=0.35,ymax=0.55,ytick={0.35,0.40,0.45,0.50,0.55},
  ymajorgrids,major grid style={black!12},tick label style={font=\scriptsize},label style={font=\scriptsize},
  title={А. Извънизвадкова $r(t)$ по часове (0--72\,ч)},title style={font=\footnotesize\bfseries}]
\addplot[thick,refcolor,mark=*,mark size=1.8pt] coordinates {(0,0.4745) (3,0.4666) (6,0.4872) (12,0.4978) (24,0.4798) (48,0.4907) (72,0.4725)};
\addplot[dashed,varcolor,thick] coordinates {(0,0.5060) (72,0.5060)};
\node[font=\scriptsize,varcolor,anchor=south east] at (axis cs:72,0.509) {Средно $r=0{,}5060$};
\end{axis}
\end{tikzpicture}\hfill
\begin{tikzpicture}
\begin{axis}[ybar,bar width=9pt,width=68mm,height=38mm,symbolic x coords={Top-25,Top-50,Top-100,База},xtick=data,
  ymin=50,ymax=105,ylabel={Реална репресия (\%)},enlarge x limits=0.2,
  ymajorgrids,major grid style={black!12},tick label style={font=\scriptsize},label style={font=\scriptsize},
  title={Б. Точност на лабораторния подбор},title style={font=\footnotesize\bfseries}]
\addplot[fill=refcolor,draw=none] coordinates {(Top-25,100.0) (Top-50,100.0) (Top-100,99.0) (База,69.1)};
\end{axis}
\end{tikzpicture}
\caption{Емпирично доказателство за Етап 1 върху 1895 строго невиждани човешки редакции (118 локуса, \code{GSE188264}). Вляво: извънизвадкова корелация на Пиърсън във всяка времева точка (0--72\,ч). Вдясно: реална лабораторна репресия за Top-25 (100\%), Top-50 (100\%) и Top-100 (99\%) обратно подбрани редакции ($p=0{,}0005$) спрямо случайната база 69,1\%.}\label{fig:v2f_hours}
\end{figure}

\section{Етапи 2--4: Двукритериално обратно спасяване на виртуална клетка и криптографско доказателство}

\subsection{Парадоксът на пасивния репортер и Парето-оптималният обратен дизайн}

Ключов недостатък на чисто емпиричния MPRA скрининг е \emph{парадоксът на пасивния репортер}. При реалните човешки енхансери \code{chr18:56475717-56475888} (\code{165:G>T}) и \code{chr1:175099416-175099587} (\code{149:G>T}) единични нуклеотидни замени променят пасивната промоторна заетост с едва $-0{,}0119$ и $-0{,}0120$, но в спрегнатата ГРМ на вентралната неврална тръба претърпяват \textbf{78,2$\times$} и \textbf{77,8$\times$} нелинейно усилване — сриват \emph{Nkx2.2} с $-0{,}9317$, превключват клетъчната идентичност към \emph{Olig2}, свиват сегментите на невритите ($34\to 32$) и разтягат общата триизмерна дължина в CX3D ($351{,}14\to 373{,}82$--$373{,}83$).

За да решим най-трудната постановка на обратното регулаторно инженерство, заключваме първичната мутация \textbf{непокътната} и оптимизираме дисталните еднобазови компенсаторни редакции едновременно по \textbf{два конкуриращи се критерия}:
\begin{enumerate}
\item \textbf{Механистична вярност на развитието:} Минимизиране на отклонението в установилите се концентрации на четирите гена, максимизиране на устойчивостта в 24-членния ансамбъл за параметрична неопределеност и възстановяване на триизмерната морфология в CX3D и спайковата проводимост (LIF).
\item \textbf{Глобална 532-ТФ нецелева безопасност:} Едновременна оценка на всеки двойномутантен кандидат чрез \code{epigenesis-biophysical-temporal-v2} за всички 532 човешки ТФ (0--72\,ч), за да се минимизира пиковото нецелево регулаторно отклонение $|\Delta\text{MPRA}_{0\text{--}72\text{ч}}|$.
\end{enumerate}

\subsection{Машинно верифицирани Парето решения и защита срещу подмяна}

Таблица~\ref{tab:pareto} обобщава двукритериалните решения, проверени на Python и самостоятелен Rust валидатор (\code{epigenesis-validate-rs}) с GA4GH VRS~2.x и RO-Crate~1.1 / COMBINE OMEX архиви:
\begin{itemize}
\item \textbf{\code{chr18:56475717-56475888} (заключена \code{165:G>T}, \code{VA.Tb8SIdM4\ldots}):} Epigenesis открива две работещи дистални спасявания (\code{163:T>C} и \code{166:G>T}). Парето-оптимизацията избира \textbf{\code{163:T>C}} ($+2$~бази, \code{VA.DXH7NbxW\ldots}), която строго доминира \code{166:G>T} и по \emph{двата} критерия: максималната генна грешка спада $12\,400\times$ (от $1{,}23\times10^{-3}$ до \textbf{$9{,}89\times10^{-8}$}), глобалното пиково 532-ТФ отклонение за 0--72\,ч спада с \textbf{23,8\%} ($0{,}1497\to 0{,}1140$; средно $-0{,}0700$ срещу $-0{,}1130$), триизмерната невритна архитектура в CX3D се възстановява на \textbf{100,00\%} ($351{,}14\to 373{,}82\to 351{,}14$; $34\to 32\to 34$ сегмента, 4 клетки, 4 синапса), а 100,0\% от ансамбловите проби възстановяват дивия тип съдба (\code{bundle\_sha256} \code{f20058a0ccd9\ldots}, \code{omex\_sha256} \code{9ceccc7a472e\ldots}).
\item \textbf{\code{chr1:175099416-175099587} (заключена \code{149:G>T}, \code{VA.Mjx89n4V\ldots}):} Компилираната дистална замяна \textbf{\code{147:A>T}} ($+2$~бази, \code{VA.9bUzqzHX\ldots}) понижава глобалното 532-ТФ пиково отклонение с \textbf{29,0\%} ($0{,}1712\to 0{,}1216$), възстановява \textbf{91,13\%} от триизмерната дължина в CX3D ($351{,}14\to 373{,}83\to 349{,}13$; $34\to 32\to 34$ сегмента) и възстановява дивия тип в \textbf{95,8\%} от пробите на ансамбъла (\code{bundle\_sha256} \code{278bd7fbd5c9\ldots}).
\item \textbf{Причинни контроли и тест за подмяна:} Всички изпълнения дават точно нулева грешка ($0{,}0$) при връщане на алела, хроматиново екраниране и изолация на пасивния репортер. И Python, и Rust валидаторите връщат \code{valid=true}; подмяната дори на 1 байт във входния FASTA предизвиква незабавно отхвърляне (\code{tamper\_attack\_rejected = true}).
\end{itemize}

\begin{table}[tb]
\scriptsize\setlength{\tabcolsep}{4pt}
\caption{Доказателство за Етапи 2--4: Парето-оптимални двукритериални обратни спасявания на виртуална клетка при заключена първична мутация (Python + Rust).}\label{tab:pareto}
\begin{tabular}{@{}L{30mm}L{19mm}r r r r L{21mm}@{}}
\toprule
Човешки енхансер & Парето спасяване & Макс.\ $\Delta$ & 532-ТФ спад & 3D CX3D & Анс.\ \% & Py / Rust / Атака \\
\midrule
\code{chr18:56475717-\allowbreak 56475888} (\code{165:G>T}, $78{,}2\times$) & \textbf{\code{163:T>C}} ($+2$\,б.) & \textbf{$9{,}89\times10^{-8}$} & \textbf{$-23{,}8\%$} ($0{,}150\to 0{,}114$) & \textbf{100,00\%} & \textbf{100,0\%} & Да · Да · Отхвърля \\
\code{chr1:175099416-\allowbreak 175099587} (\code{149:G>T}, $77{,}8\times$) & \textbf{\code{147:A>T}} ($+2$\,б.) & $1{,}11\times10^{-3}$ & \textbf{$-29{,}0\%$} ($0{,}171\to 0{,}122$) & \textbf{91,13\%} & \textbf{95,8\%} & Да · Да · Отхвърля \\
\bottomrule
\end{tabular}
\end{table}

\section{Честна научна равносметка: какво е решено и какво остава отворено}

Голямото предизвикателство изисква точна научна отчетност без преувеличения:

\textbf{Какво Epigenesis решава реално днес (доказателство \code{fa135804ba0b\ldots}):}
\begin{enumerate}
\item \textbf{Интерпретируемо, високоточно некодиращо V2F при човешка неврогенеза:} Надминаване на $r=0{,}5060$ (95\% ДИ $[0{,}3938;\, 0{,}6128]$) и 99--100\% точност на лабораторна репресия в Top-25/50/100 ($p=0{,}0005$) върху 1895 невиждани редакции в 118 отделени човешки локуса за милисекунди на CPU.
\item \textbf{Многомащабен Парето обратен дизайн с нецелева безопасност:} Едновременно решаване на епистатично триизмерно спасяване при непокътната първична мутация (91--100\% възстановяване) и 23,8\%--29,0\% намаляване на глобалното 532-ТФ отклонение.
\item \textbf{Междуезикова криптографска възпроизводимост:} Независимо преизпълнение на Python и Rust с GA4GH VRS~2.x, RO-Crate/OMEX и незабавно отхвърляне на еднобайтова подмяна.
\end{enumerate}

\textbf{Какво остава отворено за изчислителната биология:}
\begin{enumerate}
\item \textbf{Директен пренос между различни анализи и клетъчни типове:} Тринайсетте фиксирани теста в \code{research/} показват, че само локалните последователностни признаци остават около $r\approx 0{,}049$ при пренос между разнородни протоколи; преодоляването изисква далекодействащи хроматинови представяния, калибрирани спрямо клетъчно-специфичната експресия на ТФ.
\item \textbf{Далекодействащи 3D хроматинови бримки ($>10$\,kb) и ендогенно CRISPR калибриране:} Съвременните енхансери са 171--500 бази и ползват литературни ОДУ/CX3D единици; директното калибриране на синтезираните спасявания спрямо ендогенни CRISPR експерименти в човешки органоиди е решителната следваща стъпка.
\end{enumerate}

\section{Възпроизводимост}

Всяко число, VRS~2.x идентификатор и SHA-256 контролна сума в този доклад се възпроизвеждат детерминистично за $\approx 22$\,с чрез:
\begin{center}
\code{python3 examples/developmental-pipeline/grand\_challenge.py --output /tmp/epigenesis-gc}
\end{center}
Всички модули работят нативно в macOS SwiftUI/AppKit приложението (\code{Epigenesis-1.7.0-macOS.dmg}; нативните версии за Linux и Windows предстоят скоро\texttrademark) без браузърна обвивка и без локални мрежови портове.

\begin{thebibliography}{99}
\setlength{\itemsep}{2pt}

\bibitem[Agarwal et~al.(2025)]{agarwal2025} Agarwal V, Inoue F, Schubach M, et al. (2025). Massively parallel characterization of transcriptional regulatory elements. \textit{Nature} 639:411–420. \href{https://doi.org/10.1038/s41586-024-08430-9}{doi:10.1038/s41586-024-08430-9}

\bibitem[Ashuach et~al.(2019)]{ashuach2019} Ashuach T, Fischer DS, Kreimer A, et al. (2019). MPRAnalyze: statistical framework for massively parallel reporter assays. \textit{Genome Biology} 20:183. \href{https://doi.org/10.1186/s13059-019-1787-z}{doi:10.1186/s13059-019-1787-z}

\bibitem[Avsec et~al.(2021)]{avsec2021} Avsec Ž, Agarwal V, Visentin D, et al. (2021). Effective gene expression prediction from sequence by integrating long-range interactions. \textit{Nature Methods} 18:1196–1203. \href{https://doi.org/10.1038/s41592-021-01252-x}{doi:10.1038/s41592-021-01252-x}

\bibitem[Aygün et~al.(2021)]{aygun2021} Aygün N, Elwell AL, Liang D, et al. (2021). Brain-trait-associated variants impact cell-type-specific gene regulation during neurogenesis. \textit{American Journal of Human Genetics} 108:1647–1668. \href{https://doi.org/10.1016/j.ajhg.2021.07.011}{doi:10.1016/j.ajhg.2021.07.011}

\bibitem[Balaskas et~al.(2012)]{balaskas2012} Balaskas N, Ribeiro A, Panovska J, et al. (2012). Gene regulatory logic for reading the Sonic Hedgehog signaling gradient in the vertebrate neural tube. \textit{Cell} 148:273–284. \href{https://doi.org/10.1016/j.cell.2011.10.047}{doi:10.1016/j.cell.2011.10.047}

\bibitem[Beal et~al.(2011)]{beal2011} Beal J, Lu T, Weiss R. (2011). Automatic compilation from high-level biologically-oriented programming language to genetic regulatory networks. \textit{PLoS ONE} 6:e22490. \href{https://doi.org/10.1371/journal.pone.0022490}{doi:10.1371/journal.pone.0022490}

\bibitem[Bergmann et~al.(2014)]{bergmann2014} Bergmann FT, Adams R, Moodie S, et al. (2014). COMBINE archive and OMEX format: one file to share all information to reproduce a modeling project. \textit{BMC Bioinformatics} 15:369. \href{https://doi.org/10.1186/s12859-014-0369-z}{doi:10.1186/s12859-014-0369-z}

\bibitem[Bilitchenko et~al.(2011)]{bilitchenko2011} Bilitchenko L, Liu A, Cheung S, et al. (2011). Eugene – a domain specific language for specifying and constraining synthetic biological parts, devices, and systems. \textit{PLoS ONE} 6:e18882. \href{https://doi.org/10.1371/journal.pone.0018882}{doi:10.1371/journal.pone.0018882}

\bibitem[Cai et~al.(2009)]{cai2009} Cai Y, Lux MW, Adam L, Peccoud J. (2009). Modeling structure-function relationships in synthetic DNA sequences using attribute grammars. \textit{PLoS Computational Biology} 5:e1000529. \href{https://doi.org/10.1371/journal.pcbi.1000529}{doi:10.1371/journal.pcbi.1000529}

\bibitem[Chen et~al.(2022)]{chen2022} Chen KM, Wong AK, Troyanskaya OG, Zhou J. (2022). A sequence-based global map of regulatory activity for deciphering human genetics. \textit{Nature Genetics} 54:940–949. \href{https://doi.org/10.1038/s41588-022-01102-2}{doi:10.1038/s41588-022-01102-2}

\bibitem[Cohen et~al.(2014)]{cohen2014} Cohen M, Page KM, Perez-Carrasco R, Barnes CP, Briscoe J. (2014). A theoretical framework for the regulation of Shh morphogen-controlled gene expression. \textit{Development} 141:3868–3878. \href{https://doi.org/10.1242/dev.112573}{doi:10.1242/dev.112573}

\bibitem[Deng et~al.(2024)]{deng2024} Deng C, Whalen S, Steyert M, et al. (2024). Massively parallel characterization of regulatory elements in the developing human cortex. \textit{Science} 384:eadh0559. \href{https://doi.org/10.1126/science.adh0559}{doi:10.1126/science.adh0559}

\bibitem[Frankish et~al.(2021)]{frankish2021} Frankish A, et al. (2021). GENCODE 2021. \textit{Nucleic Acids Research} 49:D916–D923. \href{https://doi.org/10.1093/nar/gkaa1087}{doi:10.1093/nar/gkaa1087}

\bibitem[Fromer et~al.(2016)]{fromer2016} Fromer M, Roussos P, Sieberts SK, et al. (2016). Gene expression elucidates functional impact of polygenic risk for schizophrenia. \textit{Nature Neuroscience} 19:1442–1453. \href{https://doi.org/10.1038/nn.4399}{doi:10.1038/nn.4399}

\bibitem[Gneiting and Raftery(2007)]{gneiting2007} Gneiting T, Raftery AE. (2007). Strictly proper scoring rules, prediction, and estimation. \textit{Journal of the American Statistical Association} 102:359–378. \href{https://doi.org/10.1198/016214506000001437}{doi:10.1198/016214506000001437}

\bibitem[Gong et~al.(2020)]{gong2020} Gong et al. (2020). The requirement of Sox2 for the spinal cord motor neuron development of zebrafish. \textit{Frontiers in Molecular Neuroscience} 13:34. \href{https://doi.org/10.3389/fnmol.2020.00034}{doi:10.3389/fnmol.2020.00034}

\bibitem[GTEx Consortium(2020)]{gtex2020} GTEx Consortium (2020). The GTEx Consortium atlas of genetic regulatory effects across human tissues. \textit{Science} 369:1318–1330. \href{https://doi.org/10.1126/science.aaz1776}{doi:10.1126/science.aaz1776}

\bibitem[Guo et~al.(2023)]{guo2023} Guo et al. (2023). Integrative analyses highlight functional regulatory variants associated with neuropsychiatric diseases. \textit{Nature Genetics} 55:1876–1891. \href{https://doi.org/10.1038/s41588-023-01533-5}{doi:10.1038/s41588-023-01533-5}

\bibitem[Hallikas et~al.(2006)]{hallikas2006} Hallikas O, Palin K, Sinjushina N, et al. (2006). Genome-wide prediction of mammalian enhancers based on analysis of transcription-factor binding affinity. \textit{Cell} 124:47–59. \href{https://doi.org/10.1016/j.cell.2005.10.042}{doi:10.1016/j.cell.2005.10.042}

\bibitem[Herrera-Delgado et~al.(2020)]{herreradelgado2020} Herrera-Delgado E, Briscoe J, Sollich P. (2020). Tractable nonlinear memory functions as a tool to capture and explain dynamical behaviors. \textit{Physical Review Research} 2:043069. \href{https://doi.org/10.1103/PhysRevResearch.2.043069}{doi:10.1103/PhysRevResearch.2.043069}

\bibitem[Huber(1964)]{huber1964} Huber PJ. (1964). Robust estimation of a location parameter. \textit{Annals of Mathematical Statistics} 35:73–101. \href{https://doi.org/10.1214/aoms/1177703732}{doi:10.1214/aoms/1177703732}

\bibitem[Jaffe et~al.(2018)]{jaffe2018} Jaffe AE, Straub RE, Shin JH, et al. (2018). Developmental and genetic regulation of the human cortex transcriptome illuminate schizophrenia pathogenesis. \textit{Nature Neuroscience} 21:1117–1125. \href{https://doi.org/10.1038/s41593-018-0197-y}{doi:10.1038/s41593-018-0197-y}

\bibitem[Jerber et~al.(2021)]{jerber2021} Jerber J, Seaton DD, Cuomo ASE, et al. (2021). Population-scale single-cell RNA-seq profiling across dopaminergic neuron differentiation. \textit{Nature Genetics} 53:304–312. \href{https://doi.org/10.1038/s41588-021-00801-6}{doi:10.1038/s41588-021-00801-6}

\bibitem[Kang et~al.(2011)]{kang2011} Kang HJ, Kawasawa YI, Cheng F, et al. (2011). Spatio-temporal transcriptome of the human brain. \textit{Nature} 478:483–489. \href{https://doi.org/10.1038/nature10523}{doi:10.1038/nature10523}

\bibitem[Kelley et~al.(2016)]{kelley2016} Kelley DR, Snoek J, Rinn JL. (2016). Basset: learning the regulatory code of the accessible genome with deep convolutional neural networks. \textit{Genome Research} 26:990–999. \href{https://doi.org/10.1101/gr.200535.115}{doi:10.1101/gr.200535.115}

\bibitem[Kelley et~al.(2018)]{kelley2018} Kelley DR, Reshef YA, Bileschi M, et al. (2018). Sequential regulatory activity prediction across chromosomes with convolutional neural networks. \textit{Genome Research} 28:739–750. \href{https://doi.org/10.1101/gr.227819.117}{doi:10.1101/gr.227819.117}

\bibitem[Kerimov et~al.(2021)]{kerimov2021} Kerimov N, Hayhurst JD, Peikova K, et al. (2021). A compendium of uniformly processed human gene expression and splicing quantitative trait loci. \textit{Nature Genetics} 53:1290–1299. \href{https://doi.org/10.1038/s41588-021-00924-w}{doi:10.1038/s41588-021-00924-w}

\bibitem[Kircher et~al.(2019)]{kircher2019} Kircher M, Xiong C, Martin B, et al. (2019). Saturation mutagenesis of twenty disease-associated regulatory elements at single base-pair resolution. \textit{Nature Communications} 10:3583. \href{https://doi.org/10.1038/s41467-019-11526-w}{doi:10.1038/s41467-019-11526-w}

\bibitem[Kreimer et~al.(2022)]{kreimer2022} Kreimer A, Ashuach T, Inoue F, et al. (2022). Massively parallel reporter perturbation assays uncover temporal regulatory architecture during neural differentiation. \textit{Nature Communications} 13:1504. \href{https://doi.org/10.1038/s41467-022-28659-0}{doi:10.1038/s41467-022-28659-0}

\bibitem[Lee et~al.(2025)]{lee2025} Lee S, McAfee JC, Lee J, et al. (2025). Massively parallel reporter assay investigates shared genetic variants of eight psychiatric disorders. \textit{Cell} 188:1409–1424.e21. \href{https://doi.org/10.1016/j.cell.2024.12.022}{doi:10.1016/j.cell.2024.12.022}

\bibitem[Majoros et~al.(2020)]{majoros2020} Majoros WH, Kim YS, Barrera A, et al. (2020). Bayesian estimation of genetic regulatory effects in high-throughput reporter assays. \textit{Bioinformatics} 36:331–338. \href{https://doi.org/10.1093/bioinformatics/btz545}{doi:10.1093/bioinformatics/btz545}

\bibitem[Matoba et~al.(2026)]{matoba2026} Matoba N, McAfee JC, Krupa O, et al. (2026). Massively parallel assessment of gene regulatory activity at human cortical-structure-associated variants. \textit{Nature Neuroscience}. \href{https://doi.org/10.1038/s41593-026-02454-2}{doi:10.1038/s41593-026-02454-2}

\bibitem[McAfee et~al.(2023)]{mcafee2023} McAfee JC, Lee S, Lee J, et al. (2023). Systematic investigation of allelic regulatory activity of schizophrenia-associated common variants. \textit{Cell Genomics} 3:100404. \href{https://doi.org/10.1016/j.xgen.2023.100404}{doi:10.1016/j.xgen.2023.100404}

\bibitem[Melnikov et~al.(2012)]{melnikov2012} Melnikov A, Murugan A, Zhang X, et al. (2012). Systematic dissection and optimization of inducible enhancers in human cells using a massively parallel reporter assay. \textit{Nature Biotechnology} 30:271–277. \href{https://doi.org/10.1038/nbt.2137}{doi:10.1038/nbt.2137}

\bibitem[Miller et~al.(2014)]{miller2014} Miller JA, Ding SL, Sunkin SM, et al. (2014). Transcriptional landscape of the prenatal human brain. \textit{Nature} 508:199–206. \href{https://doi.org/10.1038/nature13185}{doi:10.1038/nature13185}

\bibitem[Myint et~al.(2019)]{myint2019} Myint L, Avramopoulos DG, Goff LA, Hansen KD. (2019). Linear models enable powerful differential activity analysis in massively parallel reporter assays. \textit{BMC Genomics} 20:209. \href{https://doi.org/10.1186/s12864-019-5556-x}{doi:10.1186/s12864-019-5556-x}

\bibitem[Ng et~al.(2017)]{ng2017} Ng B, White CC, Klein HU, et al. (2017). An xQTL map integrates the genetic architecture of the human brain's transcriptome and epigenome. \textit{Nature Neuroscience} 20:1418–1426. \href{https://doi.org/10.1038/nn.4632}{doi:10.1038/nn.4632}

\bibitem[Liu et~al.(2025)]{liu2025} Liu et al. (2025). Predicting the DNA binding specificity of transcription factor mutants using family-level biophysically interpretable machine learning. \textit{Nucleic Acids Research} 53:gkaf831. \href{https://doi.org/10.1093/nar/gkaf831}{doi:10.1093/nar/gkaf831}

\bibitem[Nielsen et~al.(2016)]{nielsen2016} Nielsen AAK, Der BS, Shin J, et al. (2016). Genetic circuit design automation. \textit{Science} 352:aac7341. \href{https://doi.org/10.1126/science.aac7341}{doi:10.1126/science.aac7341}

\bibitem[NIST(2015)]{nist2015} NIST (2015). \textit{Secure Hash Standard (SHS)}, FIPS 180-4. \url{https://csrc.nist.gov/pubs/fips/180-4/upd1/final}

\bibitem[Peterson et~al.(2012)]{peterson2012} Peterson KA, Nishi Y, Ma W, et al. (2012). Neural-specific Sox2 input and differential Gli-binding affinity provide context and positional information in Shh-directed neural patterning. \textit{Genes \& Development} 26:2802–2816. \href{https://doi.org/10.1101/gad.207142.112}{doi:10.1101/gad.207142.112}

\bibitem[Pollard et~al.(2010)]{pollard2010} Pollard KS, Hubisz MJ, Rosenbloom KR, Siepel A. (2010). Detection of nonneutral substitution rates on mammalian phylogenies. \textit{Genome Research} 20:110–121. \href{https://doi.org/10.1101/gr.097857.109}{doi:10.1101/gr.097857.109}

\bibitem[Ramasamy et~al.(2014)]{ramasamy2014} Ramasamy A, Trabzuni D, Guelfi S, et al. (2014). Genetic variability in the regulation of gene expression in ten regions of the human brain. \textit{Nature Neuroscience} 17:1418–1428. \href{https://doi.org/10.1038/nn.3801}{doi:10.1038/nn.3801}

\bibitem[Rauluseviciute et~al.(2024)]{rauluseviciute2024} Rauluseviciute I, Riudavets-Puig R, Blanc-Mathieu R, et al. (2024). JASPAR 2024: 20th anniversary of the open-access database of transcription factor binding profiles. \textit{Nucleic Acids Research} 52:D174–D182. \href{https://doi.org/10.1093/nar/gkad1059}{doi:10.1093/nar/gkad1059}

\bibitem[Roadmap Epigenomics Consortium(2015)]{roadmap2015} Roadmap Epigenomics Consortium, Kundaje A, Meuleman W, et al. (2015). Integrative analysis of 111 reference human epigenomes. \textit{Nature} 518:317–330. \href{https://doi.org/10.1038/nature14248}{doi:10.1038/nature14248}

\bibitem[Rummel et~al.(2023)]{rummel2023} Rummel CK, Gagliardi M, Ahmad R, et al. (2023). Massively parallel functional dissection of schizophrenia-associated noncoding genetic variants. \textit{Cell} 186:5165–5182.e33. \href{https://doi.org/10.1016/j.cell.2023.09.015}{doi:10.1016/j.cell.2023.09.015}

\bibitem[Rundgren et~al.(2020)]{rundgren2020} Rundgren A, Jordan B, Erdtman S. (2020). \textit{JSON Canonicalization Scheme (JCS)}, RFC 8785. \url{https://www.rfc-editor.org/rfc/rfc8785}

\bibitem[Salomon et~al.(2026)]{salomon2026} Salomon K, Deng C, Dash PM, et al. (2026). Massively parallel characterization and predictive modelling of neuronal regulatory variation. \textit{bioRxiv} (preprint, posted 17 July 2026). \href{https://doi.org/10.64898/2026.07.16.738760}{doi:10.64898/2026.07.16.738760}

\bibitem[Schwartzentruber et~al.(2018)]{schwartzentruber2018} Schwartzentruber J, Foskolou S, Kilpinen H, et al. (2018). Molecular and functional variation in iPSC-derived sensory neurons. \textit{Nature Genetics} 50:54–61. \href{https://doi.org/10.1038/s41588-017-0005-8}{doi:10.1038/s41588-017-0005-8}

\bibitem[Seitz et~al.(2024)]{seitz2024} Seitz EE, McCandlish DM, Kinney JB, Koo PK. (2024). Interpreting cis-regulatory mechanisms from genomic deep neural networks using surrogate models. \textit{Nature Machine Intelligence} 6:701–713. \href{https://doi.org/10.1038/s42256-024-00851-5}{doi:10.1038/s42256-024-00851-5}

\bibitem[Siraj et~al.(2026)]{siraj2026} Siraj L, et al. (2026). Functional dissection of complex trait variants at single-nucleotide resolution. \textit{Nature}. \href{https://doi.org/10.1038/s41586-026-10121-6}{doi:10.1038/s41586-026-10121-6}

\bibitem[Tareen et~al.(2022)]{tareen2022} Tareen A, Kooshkbaghi M, Posfai A, et al. (2022). MAVE-NN: learning genotype-phenotype maps from multiplex assays of variant effect. \textit{Genome Biology} 23:98. \href{https://doi.org/10.1186/s13059-022-02661-7}{doi:10.1186/s13059-022-02661-7}

\bibitem[Tewhey et~al.(2016)]{tewhey2016} Tewhey R, Kotliar D, Park DS, et al. (2016). Direct identification of hundreds of expression-modulating variants using a multiplexed reporter assay. \textit{Cell} 165:1519–1529. \href{https://doi.org/10.1016/j.cell.2016.04.027}{doi:10.1016/j.cell.2016.04.027}

\bibitem[Traets et~al.(2021)]{traets2021} Traets JJ, van der Burght SN, Rademakers S, et al. (2021). Mechanism of life-long maintenance of neuron identity despite molecular fluctuations. \textit{eLife} 10:e66955. \href{https://doi.org/10.7554/eLife.66955}{doi:10.7554/eLife.66955}

\bibitem[Vaishnav et~al.(2022)]{vaishnav2022} Vaishnav ED, de Boer CG, Molinet J, et al. (2022). The evolution, evolvability and engineering of gene regulatory DNA. \textit{Nature} 603:455–463. \href{https://doi.org/10.1038/s41586-022-04506-6}{doi:10.1038/s41586-022-04506-6}

\bibitem[Walker et~al.(2019)]{walker2019} Walker RL, Ramaswami G, Hartl C, et al. (2019). Genetic control of expression and splicing in developing human brain informs disease mechanisms. \textit{Cell} 179:750–771.e22. \href{https://doi.org/10.1016/j.cell.2019.09.021}{doi:10.1016/j.cell.2019.09.021}

\bibitem[Wang et~al.(2020)]{wang2020} Wang G, Sarkar A, Carbonetto P, Stephens M. (2020). A simple new approach to variable selection in regression, with application to genetic fine mapping. \textit{Journal of the Royal Statistical Society Series B} 82:1273–1300. \href{https://doi.org/10.1111/rssb.12388}{doi:10.1111/rssb.12388}

\bibitem[Weiss et~al.(2021)]{weiss2021} Weiss CV, Harshman L, Inoue F, et al. (2021). The cis-regulatory effects of modern human-specific variants. \textit{eLife} 10:e63713. \href{https://doi.org/10.7554/eLife.63713}{doi:10.7554/eLife.63713}

\bibitem[Wilson et~al.(2016)]{wilson2016} Wilson ML, et al. (2016). Genotype Specification Language. \textit{ACS Synthetic Biology} 5:471–478. \href{https://doi.org/10.1021/acssynbio.5b00194}{doi:10.1021/acssynbio.5b00194}

\bibitem[Young et~al.(2021)]{young2021} Young AMH, Kumasaka N, Calvert F, et al. (2021). A map of transcriptional heterogeneity and regulatory variation in human microglia. \textit{Nature Genetics} 53:861–868. \href{https://doi.org/10.1038/s41588-021-00875-2}{doi:10.1038/s41588-021-00875-2}

\bibitem[Zhou and Troyanskaya(2015)]{zhou2015} Zhou J, Troyanskaya OG. (2015). Predicting effects of noncoding variants with deep learning–based sequence model. \textit{Nature Methods} 12:931–934. \href{https://doi.org/10.1038/nmeth.3547}{doi:10.1038/nmeth.3547}

\bibitem[Zubler and Douglas(2009)]{zubler2009} Zubler F, Douglas R. (2009). A framework for modeling the growth and development of neurons and networks. \textit{Frontiers in Computational Neuroscience} 3:25. \href{https://doi.org/10.3389/neuro.10.025.2009}{doi:10.3389/neuro.10.025.2009}

\bibitem[Zubler et~al.(2011)]{zubler2011} Zubler F, Hauri A, Pfister S, et al. (2011). An instruction language for self-construction in the context of neural networks. \textit{Frontiers in Computational Neuroscience} 5:57. \href{https://doi.org/10.3389/fncom.2011.00057}{doi:10.3389/fncom.2011.00057}

\bibitem[Zubler et~al.(2013)]{zubler2013} Zubler F, Hauri A, Pfister S, et al. (2013). Simulating cortical development as a self constructing process: a novel multi-scale approach combining molecular and physical aspects. \textit{PLoS Computational Biology} 9:e1003173. \href{https://doi.org/10.1371/journal.pcbi.1003173}{doi:10.1371/journal.pcbi.1003173}

\end{thebibliography}


\end{document}
